%% ma: L12-B08-0126
%% lop: 12
%% chuong: 2
%% bai: 8
%% dang: 12.8.5
%% loai: TN4
%% muc: TH
%% dap_an: "A"
%% nguon: "CĐ2.3.pdf giáo viên gửi 27/10/2026 (M07, Lê Minh Khôi – Nguyễn Anh Khôi, Chương II Chủ đề 3 Biểu thức tọa độ của các phép toán vectơ), Câu 41"
%% kien_thuc: "Bài 8 (ba điểm thẳng hàng, khoảng cách)"
%% trang_thai: chinh-thuc
\begin{ex}
Trong không gian $Oxyz$, cho hai điểm $A(-2;3;1)$ và $B(5;6;2)$. Đường thẳng $AB$ cắt mặt phẳng $(Oxz)$ tại điểm $M$. Tính giá trị của biểu thức $K=\dfrac{AM}{BM}$.
\choice{\True $K=0{,}5$}{$K=0{,}25$}{$K=1$}{$K=1{,}5$}
\loigiai{$M(x;0;z)$, $\overrightarrow{AM}=t\overrightarrow{AB}=t(7;3;1)$, tung độ: $3+3t=0\Rightarrow t=-1$, $M(-9;0;0)$.

$AM=AB$, $BM=2AB$ nên $K=0{,}5$.}
\end{ex}
